\subsection{Dataset and snapshot sampling}

The analysis-ready corpus contained 27 dated dashboard snapshots spanning 1 January 2024 to 29 September 2026. There were 9 dates in each year. The 2024--2025 dates were manifest-defined candidates retained from the B0.9A/B0.9C range design, with year boundaries, month boundaries, selected dated states, and year-validation dates. The 2026 dates were the validated B0.8A/B pilot set. All 27 manifest rows were represented in the historical ledger; 16 were recorded as retrieved-valid and 11 as reused-verified artifacts, with no documented acquisition failure substitution. The available records do not establish whether every selection decision preceded viewing all possible outcomes, so the sample is described as predefined in the manifest, not as outcome-blinded or representative (Supplementary Table S3).

This design supplies an auditable set of historical states rather than a claim of continuous daily coverage. It also distinguishes newly retrieved responses from previously verified artifacts, which is relevant when a later user attempts to reproduce the archive or interpret differences between dates.

\subsection{Operational reconstruction and provenance}

All 27 archived states had recorded submitted dates, HTTP 200 responses, displayed dates matching the requested dates, the reference structural compatibility signature (\texttt{84efb756f47cbb24e1ad240bc2507576ead9cab3951e379ac0e618e7bb94a83d}), response metadata, and SHA-256 records. The signature summarizes normalized form structure, table headers, and selected chart-container IDs; it is not a semantic claim. The frozen processed outputs represented 11 data families across the 27 snapshots, yielding 297 family-level rows. Four source-labelled headline fields were retained per snapshot, yielding 108 field values. All 432 processed observations retained provenance fields. For the three B1.7 spot-validated snapshots---2024-01-01, 2025-09-29, and 2026-09-29---18 checks of date controls, displayed dates, summary-strip labels, and headline values all matched. These checks demonstrate reproducible reconstruction within the archived corpus; they do not demonstrate complete upstream surveillance coverage (Supplementary Tables S1--S3).

\subsection{Semantic and category findings}

The 11-row anomaly register shows concrete examples of why raw-label preservation was necessary. In the age-by-sex family, \texttt{0-10} overlaps with \texttt{0-5} and \texttt{6-10}; \texttt{06-10} and \texttt{6-10} are format variants that were not shown to be safely mergeable; and \texttt{42309} and \texttt{46301} are numeric-looking labels without an established age meaning. In the division family, \texttt{Raishahi} was preserved separately from \texttt{Rajshahi} rather than silently corrected. Geographic examples included \texttt{Dhaka Division (Out of CC)}, whose residence/facility/administrative meaning remained unresolved, and the raw \texttt{DNCC} and \texttt{DSCC} City Corporation labels. These labels were not converted into rural, urban-population, or patient-residence categories. Age-by-sex and death-age-by-sex families were unavailable in the frozen processed table for all 27 snapshots. The complete register and treatment decisions are in Table-Category-Anomalies and the B0.11 semantic maps.

\subsection{Internal reconciliation findings}

The nine reconciliation tests show different forms of restricted comparability. The weekly-case versus cumulative-case comparison was retained as a non-identity audit comparison, not an arithmetic equality. Age-by-sex component tests remained unresolved because raw labels and table-to-family universes were not frozen. Within/outside City Corporation and division subtotal comparisons were not comparable because geographic attribution and compatible denominators were not established. Daily, EPI-week, and monthly accumulation tests remained unresolved because typed series and compatible totals were unavailable. Weekly and cumulative death quantities were not summed because they represent different reporting windows. The B0.11 classifications therefore identify missing categories, semantic mismatch, different time windows, unknown compatibility, or dashboard-raw reconciliation limitations; they do not call these observations biological or system failures (Table-Reconciliation-Audit).

\subsection{Revision and stability findings}

The revision audit contains four entries---daily, EPI-week, monthly, and annual---and each is \texttt{UNRESOLVED}. The frozen processed parent contains family-level observations but not typed repeated chart-point values linked to a stable observation identity across captures. Establishing retrospective revision would have required repeated typed values for the same historical point, with stable time keys and comparable field definitions. Because that evidence was absent, the study did not demonstrate that revisions occurred, that the dashboard was stable, or that a revision policy existed (Table-Revision-Stability-Audit).

\subsection{Restricted descriptive observations}

The three predefined illustrative snapshots make the source-labelled headline structure visible. On 2024-01-01, the weekly and cumulative case fields were both 82 and the corresponding death fields were both 0. On 2025-09-29, the weekly case and death fields were 1,579 and 7, while cumulative case and death fields were 46,783 and 195. On 2026-09-29, the corresponding values were 6,136 and 20, and 77,672 and 239. Equality of weekly and cumulative values in an individual snapshot does not establish semantic equivalence between the fields. All three snapshots passed the 18-check B1.7 spot-validation. These are source-labelled dashboard fields, not population incidence, risk, or community infection estimates; 2026 remains incomplete (Table-Illustrative-Headline-Snapshots).
